%----------------------------------------------------------------------------------------
%	PACKAGES AND OTHER DOCUMENT CONFIGURATIONS
%----------------------------------------------------------------------------------------

\documentclass[10pt, a4paper, oneside]{article} % Paper size, default font size and one-sided paper

\usepackage{nomencl}
\usepackage{hyperref}
\usepackage{setspace}
\usepackage{fancyhdr}
\usepackage[utf8]{inputenc}
\usepackage{graphicx}
\usepackage{color}
\usepackage{epstopdf}
\usepackage[final]{pdfpages}
\usepackage[utf8]{inputenc}
\usepackage{float}

\usepackage[margin=1in]{geometry}
\usepackage[square, numbers, comma, sort&compress]{natbib} 
%\makenomenclature
\renewcommand{\nomname}{Time Zones}
\newcommand{\solidareit}{\textsc{s}olidare-\textsc{it} }
\title{\ttitle} % Defines the thesis title - don't touch this

\begin{document}
%\frontmatter % Use roman page numbering style (i, ii, iii, iv...) for the pre-content pages

\setstretch{1.2} % Line spacing of 1.3

% Define the page headers using the FancyHdr package and set up for one-sided printing
\fancyhead{} % Clears all page headers and footers
\rhead{\thepage} % Sets the right side header to show the page number
\lhead{} % Clears the left side page header

\pagestyle{fancy} % Finally, use the "fancy" page style to implement the FancyHdr headers

\newcommand{\HRule}{\rule{\linewidth}{0.5mm}} % New command to make the lines in the title page

%----------------------------------------------------------------------------------------
%	TITLE PAGE
%----------------------------------------------------------------------------------------


\definecolor{darkblue}{rgb}{0.1,0.3,0.7}
\definecolor{red}{rgb}{1.0,0.2,0.2}
\definecolor{darkgreen}{rgb}{0.2,0.5,0.2}


\begin{titlepage}

\begin{tabular}{cc}

\begin{minipage}{0.49\textwidth}
\begin{flushleft}
\includegraphics[scale=0.1]{./Figures/logoingisbleu.jpg} % University/department logo - uncomment to place it
\end{flushleft}
\end{minipage}

&
 \begin{minipage}{0.42\textwidth}
\begin{flushright}
\includegraphics[scale=0.5]{./Figures/epl.jpg} % University/department logo - uncomment to place it
\end{flushright}
\end{minipage}
\end{tabular} 



\begin{center}
\vspace{13em}
\textsc{\LARGE lingi2263 : Computational Linguistics }\\[2cm] % University name

 \vspace{1em}
\HRule \\[0.5cm] % Horizontal line
{\huge  Group 2 : Project 2}\\[0.35cm] % Thesis title
\HRule \\[1.5cm] % Horizontal line
 

\begin{tabular}{ccc}
\begin{minipage}{0.55\textwidth}
\begin{flushleft} \large
\emph{Authors:}\\{
Crochelet Martin (2236-10-00)\\
Baugnies Benjamin (6020-10-00)}
\end{flushleft}
\end{minipage} & \begin{minipage}{0.41\textwidth}
\centering
\begin{flushright} \large
\emph{Professor:}\\{
Pierre Dupont\\
Cédric Fairon\\
}
\end{flushright}
\end{minipage}\\[3cm] \\ 
\end{tabular} 

\vspace{4em}


 \begin{center}
{\large 2013 - 2014}\\[4cm] % Date 
 \end{center}


\vfill
\end{center}

\end{titlepage}

%----------------------------------------------------------------------------------------
%	LIST OF CONTENTS/FIGURES/TABLES PAGES
%----------------------------------------------------------------------------------------

\pagestyle{fancy} % The page style headers have been "empty" all this time, now use the "fancy" headers as defined before to bring them back

%\lhead{\vspace{1em}\emph{Contents}} % Set the left side page header to "Contents"
%\tableofcontents % Write out the Table of Contents

%\lhead{\emph{List of Figures}} % Set the left side page header to "List of Figures"
%\listoffigures % Write out the List of Figures

%\lhead{\emph{List of Tables}} % Set the left side page header to "List of Tables"
%\listoftables % Write out the List of Tables

%----------------------------------------------------------------------------------------
%	ABBREVIATIONS
%----------------------------------------------------------------------------------------
%
%\clearpage % Start a new page
%
%\setstretch{1.5} % Set the line spacing to 1.5, this makes the following tables easier to read
%
%\lhead{\emph{Abbreviations}} % Set the left side page header to "Abbreviations"
%\listofsymbols{ll} % Include a list of Abbreviations (a table of two columns)
%{
%\textbf{ANT1} & \textbf{A}\textbf{N}\textbf{T}agoniste \textbf{1} \\
%}



%UTC\nomenclature{UTC}{Coordinated Universal Time} is 3 hours behind ADT\nomenclature{ADT}{Atlantic Daylight Time} and 10 hours ahead of EST\nomenclature{EST}{Eastern Standard Time}.
%\printnomenclature
%\small\hfill Created by http://texblog.org


%----------------------------------------------------------------------------------------
%	RAPPORT CONTENT - CHAPTERS
%----------------------------------------------------------------------------------------

\section{Tags and words statistics}
The first part of the assignment consisted of creating a Lexicon based on the training corpus, and extracting some general information about the words and tags it contains. Our lexicon is limited to the 5000 most common words. Other words are replaced by the \texttt{<UNK>} token. Among these words, we extracted the 10 most common words. They are summarized along with their frequency in table \ref{words}.
\begin{table}[!h]
\centering
\begin{tabular}{ | c | c | c | }
\hline
Rank & Word & \# occurences \\ \hline
1 & THE & 59466 \\
2 & , & 49639 \\
3 & . & 41859 \\
4 & OF & 31067 \\
5 & AND & 24709 \\
6 & TO & 22190 \\
7 & A & 19718 \\
8 & IN & 18230 \\
9 & THAT & 8974 \\
10 & IS &8623 \\ \hline
\end{tabular}
\caption{Most common words}
\label{words}
\end{table}

While parsing the texts, we found 187 different tags. The 10 most and least common of these can be found below \ref{tags}.
\begin{table}[!h]
\centering
\begin{tabular}{ | c | c | c | | c | c | c | }
\hline
& Most Common & & & Least Common & \\ \hline
Rank & Tag & \# occurences  & Rank & Tag & \# occurences\\ \hline
1 & NN & 143023 & -1 & JJR+CS & 1\\
2 & IN & 104385 & -2 & JJ\$ & 1\\
3 & AT & 84266 & -3 & NN+HVD & 1\\
4 & JJ & 58365 & -4 & RBR+CS & 1\\
5 & . & 51997 & -5 & IN+IN & 1\\
6 & , & 49641 & -6 & NR+MD & 1\\
7 & NNS & 49412 & -7 & IN+NP & 1\\
8 & NP & 32992 & -8 & WRB+BER & 1\\
9 & CC & 32496 & -9 & WRB+MD & 1\\
10 & RB & 31147 & -10 & NP+MD & 1\\ \hline
\end{tabular}
\caption{Most and least common tags}
\label{tags}
\end{table}

We can see that the most rarely seen tags are almost all compound tags. Moreover, the list of least used tags is actually not very precise since a lot more than 10 tags are used only once. The presented list here is just a chosen subset of those tags but we could have chosen other tags. 

\section{Tags and words statistics}
\section{HMM Tagger}

This last part has been more difficult than foreseen: indeed, we have been able to compute the emission and transition matrix from the training set however, computing the whole set of possible states for a given segment has been quite a challenge. In the joined code, you can see that the whole HMM POS-tagger has been implemented less this challenging part. We expect the tagger to be a lot more slower than the baseline one since it has to access the matrix for each possible state which is a lot more complex than the normal, baseline, version. Moreover, our choice of smoothing is a little bit unorthodox since we simple use a default epsilon value; a better implementation would follow the formula given at slide 18 of the fifth lecture. 

About the evaluation of the tagger, we have obviously not been able to execute it however, the implementation is given in the python file compare.py. 

%\mainmatter % Begin numeric (1,2,3...) page numbering

\pagestyle{fancy} % Return the page headers back to the "fancy" style

% Include the chapters of the thesis as separate files from the Chapters folder
% Uncomment the lines as you write the chapters
\pagebreak
%\input{./Chapters/introduction.tex}
%\input{./Chapters/high_level_view_of_the_website.tex}
%\pagebreak
%\input{./Chapters/functional_requirements.tex}
%\pagebreak
%\input{./Chapters/non_functional_requirements.tex}
%\pagebreak
%\input{./Chapters/clarification_questions_and_suggestions.tex}
%\pagebreak
%\input{./Chapters/selected_extensions.tex}
%\input{./Chapters/conclusion.tex}
%\input{./Chapters/Introduction.tex}
%\input{./Chapters/ContexteGlobal.tex}
%\input{./Chapters/Localite.tex}
%\input{./Chapters/abreprob.tex}
%\input{./Chapters/Etatdelart.tex}
%\input{./Chapters/Recommandations.tex}
%\input{./Chapters/ccl.tex}




%----------------------------------------------------------------------------------------
%	BIBLIOGRAPHY
%----------------------------------------------------------------------------------------
%\newpage
%\section{Bibliography}
%\label{Bibliography}
%
%\lhead{\vspace{1em}\emph{Bibliography}} % Change the page header to say "Bibliography"
%
%\bibliographystyle{unsrtnat} % Use the "unsrtnat" BibTeX style for formatting the Bibliography
%
%\bibliography{Bibliography} % The references (bibliography) information are stored in the file named "Bibliography.bib"

%----------------------------------------------------------------------------------------
%	THESIS CONTENT - APPENDICES
%----------------------------------------------------------------------------------------

%\addtocontents{toc}{\vspace{2em}} % Add a gap in the Contents, for aesthetics
%
%\appendix % Cue to tell LaTeX that the following 'chapters' are Appendices
%
%% Include the appendices of the thesis as separate files from the Appendices folder
%% Uncomment the lines as you write the Appendices
%
%\input{./Appendices/resultats}
%\input{./Appendices/finances}
%\input{./Appendices/methodo}
%
%\addtocontents{toc}{\vspace{2em}} % Add a gap in the Contents, for aesthetics
%
%%\backmatter


\end{document}  